\documentclass[10pt,a4paper]{amsart}
\usepackage[margin=19mm]{geometry}
\usepackage{amsmath,amssymb,mathtools}
\usepackage[scr=rsfs,cal=euler]{mathalfa}
\usepackage{xfrac}
\usepackage[unicode,hidelinks,bookmarks=false]{hyperref}
\newcommand{\ie}{i.e.\ }
\newcommand{\cf}{cf.\ }
\newcommand{\resp}{respectively\ }
\newcommand{\Subsection}{\subsection}
\providecommand{\nobreakdash}{\nobreak\hbox{-}}
\setlength{\emergencystretch}{2em}
\multlinegap0pt
\usepackage{xcolor}
\usepackage{braket}
\usepackage{amssymb}
\usepackage{dsfont}
\usepackage{mathtools}
\usepackage{algorithm}\usepackage{algpseudocode}
\usepackage[shortlabels]{enumitem}
\usepackage{float}
\usepackage[normalem]{ulem}
\usepackage{tikz}
\newtheorem{lemma}{Lemma}
\usepackage{cleveref}
\crefname{lemma}{Lemma}{Lemmas}
\newcommand{\CC}{\mathbb{C}}
\newcommand{\PP}{\mathbb{P}}
\DeclareMathOperator{\Res}{Res}
\newcommand{\cR}{\mathcal{R}}
\newcommand{\ff}{\ensuremath{\boldsymbol{f}}}
\newcommand{\hh}{\ensuremath{\boldsymbol{h}}}
\DeclareMathOperator{\mult}{mult}
\title{Resultant multiplicity via projective degrees\\ and applications to tensor eigenvalues}
\author{Mahmut Levent Do\u{g}an, Elias Tsigaridas, Zafeirakis Zafeirakopoulos}
\date{Source: arXiv:2609.01268v1 (CC BY 4.0)}
\begin{document}
\maketitle

\noindent\emph{Source and licence.} Resultant multiplicity via projective degrees\\ and applications to tensor eigenvalues. Version arXiv:2609.01268v1 (CC BY 4.0), distributed by arXiv under the Creative Commons Attribution 4.0 International licence, http://creativecommons.org/licenses/by/4.0/, as recorded in the arXiv licence metadata for this exact version and shown on the abstract page https://arxiv.org/abs/2609.01268v1. Full licence notice: \texttt{licenses/arxiv-2609.01268-CC-BY-4.0.txt}.\par\smallskip

\noindent\textit{Editorial scope.} Counted unit: the article's lemma lem:count (source0.tex lines 877--882). The article's complete original proof of it is delivered with it (source0.tex lines 883--910, one \texttt{proof} environment of 28 lines). No mathematics is written, repaired or re-derived: the statement and the proof are the article's own lines, in the article's own notation, and no retained line is substituted or re-flowed.  Retained uncounted as the source-local context the proof needs: lines 478--483; lines 768--777.  Nothing was omitted from any retained span, so the package carries no omission record.  No retained line contains a citation, so the wrapper carries no bibliography and no bibliography entry.  No retained line contains a figure environment, an image-inclusion command or drawing code, so no image is delivered and the package declares no asset dependency.  Editorial formatting only, in addition: an AMS \texttt{amsart} preamble that restates the article's own declarations and macros used by the retained lines, namely the environments \texttt{lemma} on separate counters, together with the standard packages those lines need.  The retained environments print in this document as Lemma 1 in order of appearance, whereas the article prints the same items with the numbering of its own full document.  The complete native source of the article as distributed in the arXiv e-print for this exact version is bundled unchanged as \texttt{sources/209090/source0.tex}.
\begin{lemma}
\label{lem:generic-line}
    For every $q\in\CC^m$ we have \[
    \operatorname{ord} F(p+tq) \geq \mult_{p}(F),
    \] and equality holds for generic $q$.
\end{lemma}

\begin{lemma}
\label{lem:coincidence-pts}
    For a general $\hh$ as above, the intersection \[
    \Gamma_{\ff}\cap \Gamma_{\hh}
    \] is contained in $G_{\ff}$, is transverse, and is therefore zero-dimensional and reduced.
    It is isomorphic to $\mathcal{C}(\ff,\hh)$ by projection onto the first coordinate.
    Moreover, \[
    C(\ff,\hh) = \sum_{\ell=0}^{n-1} d^{n-\ell-1} y_\ell.
    \]
\end{lemma}

\begin{lemma}
\label{lem:count}
We have \[
\mu(\ff) = n d^{n-1} - C(\ff,\hh).
\]
\end{lemma}

\begin{proof}
Since $L$ is a general line through $[\ff]$, its local intersection multiplicity with the resultant hypersurface at $[\ff]$ is precisely $\mu(\ff)$, see \cref{lem:generic-line}.

Let us denote by $\cR_{\PP}\subset\PP(R_d^n)$ the projectivization of the resultant hypersurface $\cR$.
After choosing $\hh$ generically, all the other intersections of $L$ with $\cR_{\PP}$ are transverse and belong to the open, dense subset of $\cR_{\PP}$ that consists of systems with a unique, simple projective zero.
Since the total degree of the resultant is \[
\deg(\Res) = n d^{n-1},
\] 
we have \[
nd^{n-1} = \mu(\ff) + \# \left((L\cap \cR_{\PP})\setminus \{[\ff]\}\right),
\] where the intersections counted in the second summand are simple.

Since $\Res(\hh)\neq 0$, the point $[\hh]$ does not belong to the resultant hypersurface, and, therefore, does not contribute to the intersection of $L$ and $\cR_{\PP}$.
Every other point of $L\setminus\{[\ff],[\hh]\}$ is of the form $[\ff-s\hh]$ for some unique, non-zero $s\in\CC^{\times}$.
For such $s$, the system $[\ff-s\hh]$ lies in $\cR_{\PP}$ if and only if \begin{equation}
    \label{eq:sol}
 (\ff - s \hh)(p)=0
\end{equation} for some $[p]\in\PP^{n-1}$. 
Since $\hh(p)\neq 0$, this implies $\ff(p)\neq 0$ and \cref{eq:sol} gives \[
\phi_{\ff}([p]) = \phi_{\hh}([p]). 
\]
Conversely, every solution of this equality determines a unique $s\in\CC^{\times}$ such that $(\ff-s\hh)(p)=0$, since we assumed that the intersections of $L$ and $\cR_{\PP}$ occur at systems with a unique, simple zero. 

For general $\hh$, these coincidence points are simple by \cref{lem:coincidence-pts}.
Hence,  \[
 \#\left((L\cap \cR_{\PP})\setminus \{[\ff]\}\right) = C(\ff,\hh).
\]
\end{proof}

\end{document}
